% ============================================================
% Paper 3: PQC for Bandwidth-Constrained Aviation Datalinks
% IEEE Transactions on Aerospace and Electronic Systems
% ============================================================
\documentclass[conference,10pt]{IEEEtran}

\usepackage{amsmath,amssymb,amsfonts}
\usepackage{algorithmic}
\usepackage{algorithm}
\usepackage{graphicx}
\usepackage{textcomp}
\usepackage{xcolor}
\usepackage{booktabs}
\usepackage{multirow}
\usepackage{hyperref}
\usepackage{cite}
\usepackage{url}
\usepackage{balance}
\usepackage{array}
\usepackage{tabularx}
\usepackage{subcaption}

\hypersetup{
  colorlinks=true,
  linkcolor=blue,
  citecolor=blue,
  urlcolor=blue
}

\begin{document}

\title{Post-Quantum Cryptography for Bandwidth-Constrained Aviation Datalinks: Aggregate Signatures and Forward-Secure Channels for ATC Communication}

\author{
\IEEEauthorblockN{Prabakaran Kannan}
\IEEEauthorblockA{Research Scholar, Department of Computer Science and Engineering\\
National Institute of Technology Puducherry\\
Email: cs23d2005@nitpy.ac.in}
}

\maketitle

% ============================================================
% ABSTRACT
% ============================================================
\begin{abstract}
Air traffic control (ATC) communication systems operate under stringent bandwidth constraints that pose a fundamental challenge to the adoption of post-quantum cryptography (PQC). Current ATC datalinks---including VHF ACARS at 2.4~kbps, HF Datalink at 1.8~kbps, and classic SATCOM at 600~bps---cannot accommodate the significantly larger signature and key sizes of lattice-based schemes such as ML-DSA and Falcon, where individual signatures range from 666~bytes to 3.3~KB. At the same time, the ``harvest now, decrypt later'' (HNDL) threat model makes quantum-resistant security an urgent requirement for safety-critical aviation messages that may retain sensitivity for decades.

This paper presents two complementary constructions that enable practical PQC deployment on bandwidth-constrained aviation datalinks. First, we introduce a Merkle tree aggregate signature scheme that compresses $N$ individual PQC signatures into a single compact aggregate proof. By computing SHA3-256 leaf hashes of individual signatures, constructing a Merkle tree, and applying Zstd compression, the scheme achieves bandwidth reductions of 80--90\% on the most constrained channels, reducing per-signature overhead to approximately 100--200~bytes. The scheme incorporates priority-aware handling where DISTRESS messages bypass aggregation to ensure zero additional latency for safety-critical communications. Second, we design a forward-secure double-ratchet channel protocol using ephemeral ML-KEM-768 key exchanges per epoch, with AES-256-GCM symmetric encryption and HKDF-SHA3-256 key derivation. The protocol supports four ratchet modes---per-message, per-epoch, per-exchange, and periodic---tailored to the operational profiles of ACARS, CPDLC, ADS-C, and LDACS channels. Key erasure is performed immediately at epoch boundaries with a one-interval retention window for message reordering. We provide detailed performance analysis across five channel profiles, demonstrating aggregate verification latency below 50~ms for batches of 64 messages and ratchet overhead compatible with real-time ATC operations. Security analysis establishes aggregate unforgeability under the collision resistance of SHA3-256 and forward secrecy under the IND-CCA2 security of ML-KEM-768.
\end{abstract}

\begin{IEEEkeywords}
Post-quantum cryptography, air traffic control, aggregate signatures, forward secrecy, ML-KEM, ML-DSA, ACARS, CPDLC, LDACS, bandwidth optimization
\end{IEEEkeywords}

% ============================================================
% I. INTRODUCTION
% ============================================================
\section{Introduction}

Aviation communication has undergone a fundamental transformation over the past three decades. The transition from analog voice-only radio to digital datalinks---Aircraft Communications Addressing and Reporting System (ACARS), Controller-Pilot Data Link Communications (CPDLC), and Automatic Dependent Surveillance-Broadcast (ADS-B)---has enabled more efficient air traffic management but has simultaneously expanded the attack surface for cyber threats~\cite{strohmeier2015security, costin2012ghost}. The safety-critical nature of these communications, where a compromised message could endanger hundreds of lives, demands the highest assurance levels for authentication and confidentiality.

The advent of cryptographically relevant quantum computers (CRQCs) poses a particularly acute threat to aviation. Unlike many communication domains where data sensitivity is transient, ATC communications carry implications for accident investigation, regulatory compliance, and national security that extend for decades. The ``harvest now, decrypt later'' (HNDL) attack model~\cite{mosca2018cybersecurity}---where adversaries record encrypted traffic today for future quantum decryption---is especially concerning for aviation, as recorded ATC transcripts could reveal flight patterns, military movements, and critical infrastructure vulnerabilities.

The National Institute of Standards and Technology (NIST) has finalized the first generation of post-quantum cryptographic standards: ML-KEM (FIPS~203)~\cite{fips203} for key encapsulation, ML-DSA (FIPS~204)~\cite{fips204} for digital signatures, and SLH-DSA (FIPS~205)~\cite{fips205} for stateless hash-based signatures. The NSA's CNSA~2.0 suite~\cite{nsa2022cnsa2} mandates PQC adoption for national security systems by 2035, a timeline that directly affects military and civil aviation infrastructure.

However, PQC adoption faces a critical obstacle in aviation: \emph{bandwidth}. The signature sizes of lattice-based schemes are orders of magnitude larger than their classical counterparts. ML-DSA-65 produces signatures of 3,309~bytes compared to 64~bytes for Ed25519, while even the more compact Falcon-512 generates 666-byte signatures. On a VHF ACARS channel operating at 2.4~kbps, transmitting a single ML-DSA-65 signature requires approximately 11~seconds---an unacceptable delay for time-critical ATC messages.

Table~\ref{tab:sig_comparison} illustrates the severity of this mismatch. The most bandwidth-constrained channels cannot feasibly carry individual PQC signatures without either unacceptable latency or significant reduction in message throughput.

\begin{table}[t]
\centering
\caption{PQC Signature Sizes vs. ATC Channel Capacities}
\label{tab:sig_comparison}
\begin{tabular}{lcc}
\toprule
\textbf{Scheme} & \textbf{Sig. Size (B)} & \textbf{TX Time @ 2.4 kbps} \\
\midrule
Ed25519 (classical) & 64 & 0.21 s \\
Falcon-512 & 666 & 2.22 s \\
ML-DSA-44 & 2,420 & 8.07 s \\
ML-DSA-65 & 3,309 & 11.03 s \\
SLH-DSA-128s & 7,856 & 26.19 s \\
\bottomrule
\end{tabular}
\end{table}

\subsection{Contributions}

This paper addresses the bandwidth-PQC conflict in aviation through two complementary constructions:

\begin{enumerate}
    \item \textbf{Merkle Tree Aggregate Signature Scheme.} We present a hash-based aggregate signature construction that compresses $N$ individual PQC signatures into a single compact proof. The scheme builds a Merkle tree over SHA3-256 hashes of $(message \| signature)$ pairs and transmits only the tree root plus truncated leaf identifiers, achieving 80--90\% bandwidth reduction. A priority-aware mechanism ensures DISTRESS messages bypass aggregation entirely.

    \item \textbf{Forward-Secure Double-Ratchet Channels.} We design a ratcheting key derivation protocol using ephemeral ML-KEM-768 key exchanges that provides forward secrecy for ATC datalinks. The protocol supports channel-specific ratchet modes (per-message, per-epoch, per-exchange) and performs immediate key erasure at epoch boundaries.
\end{enumerate}

Both constructions are implemented and evaluated against real ATC channel profiles, demonstrating practical deployment feasibility on existing aviation infrastructure. To our knowledge, this is the first work to address both signature aggregation and forward secrecy for post-quantum aviation communications under realistic bandwidth constraints.

The remainder of this paper is organized as follows. Section~\ref{sec:background} provides background on ATC communications and PQC. Section~\ref{sec:model} presents the system model and threat model. Sections~\ref{sec:aggregate} and~\ref{sec:forward} describe our two constructions. Section~\ref{sec:evaluation} presents performance evaluation results. Section~\ref{sec:security} provides security analysis. Section~\ref{sec:related} surveys related work, and Section~\ref{sec:conclusion} concludes.

% ============================================================
% II. BACKGROUND
% ============================================================
\section{Background}
\label{sec:background}

\subsection{ATC Communication Systems}

Modern ATC relies on a heterogeneous set of communication channels, each with distinct bandwidth characteristics, message formats, and operational constraints.

\textbf{ACARS} (Aircraft Communications Addressing and Reporting System) is the oldest and most widely deployed aviation datalink, providing text-based communication between aircraft and ground stations via VHF, HF, or satellite links. VHF ACARS operates at 2.4~kbps with message sizes limited to approximately 220 characters. Despite its age, ACARS remains critical for airline operational control (AOC) messages, including weight and balance data, departure clearances, and weather updates.

\textbf{CPDLC} (Controller-Pilot Data Link Communications) enables digital communication between air traffic controllers and pilots, reducing voice channel congestion. CPDLC messages follow the ICAO-defined ATN/IPS protocol stack~\cite{icao9896} and carry safety-critical instructions such as altitude clearances, route modifications, and holding instructions. Message latency requirements are stringent: routine messages must be delivered within 10~seconds, while urgent messages require sub-5-second delivery.

\textbf{ADS-B} (Automatic Dependent Surveillance-Broadcast) transmits aircraft position, velocity, and identification at 1~MHz (1~Mbps) using 112-bit messages at 2~Hz intervals~\cite{do350a}. While ADS-B has higher raw bandwidth than ACARS, individual message sizes are severely limited, and the broadcast nature creates unique security challenges~\cite{strohmeier2015security}.

\textbf{ADS-C} (Automatic Dependent Surveillance-Contract) provides surveillance data over point-to-point datalinks as part of contractual agreements between aircraft and ATC centers. ADS-C operates over ACARS or ATN networks with periodic position reports at intervals from 30~seconds to 30~minutes.

\textbf{LDACS} (L-Band Digital Aeronautical Communications System) is the next-generation terrestrial datalink currently under standardization by EUROCAE~\cite{ed154, ed228}. LDACS provides approximately 100~kbps of capacity with built-in security mechanisms, representing a significant bandwidth improvement over legacy channels~\cite{maeurer2021ldacs}.

\subsection{Current Aviation Security: ARINC 823}

The primary security specification for ACARS datalinks is ARINC~823~\cite{arinc823}, which defines the AMS (ACARS Message Security) protocol. ARINC~823 provides message authentication via HMAC-SHA-256 and optional confidentiality via AES-128-CBC, with pre-shared symmetric keys distributed through airline key management systems.

ARINC~823 has several limitations relevant to the quantum threat:
\begin{itemize}
    \item Key establishment relies on RSA or ECDH, both vulnerable to Shor's algorithm~\cite{shor1997polynomial}.
    \item No forward secrecy: compromise of the long-term key reveals all past and future session keys derived from it.
    \item No signature aggregation: each message carries its own authentication tag, consuming bandwidth individually.
    \item Limited crypto-agility: transitioning to PQC requires protocol-level changes beyond simple algorithm substitution.
\end{itemize}

\subsection{Post-Quantum Signature Size Challenge}

The fundamental tension between PQC and aviation bandwidth is quantitative. Table~\ref{tab:pqc_sizes} summarizes the signature and public key sizes for NIST-standardized PQC signature schemes at security level~2 (128-bit equivalent) and level~3 (192-bit equivalent).

\begin{table}[t]
\centering
\caption{NIST PQC Signature Scheme Sizes}
\label{tab:pqc_sizes}
\begin{tabular}{lccc}
\toprule
\textbf{Scheme} & \textbf{PK (B)} & \textbf{Sig (B)} & \textbf{Level} \\
\midrule
ML-DSA-44 & 1,312 & 2,420 & 2 \\
ML-DSA-65 & 1,952 & 3,309 & 3 \\
ML-DSA-87 & 2,592 & 4,627 & 5 \\
Falcon-512 & 897 & 666 & 1 \\
Falcon-1024 & 1,793 & 1,280 & 5 \\
SLH-DSA-128s & 32 & 7,856 & 1 \\
SLH-DSA-128f & 32 & 17,088 & 1 \\
\bottomrule
\end{tabular}
\end{table}

Even the most compact PQC signature scheme (Falcon-512 at 666~bytes) produces signatures over 10$\times$ larger than ECDSA-256 (64~bytes). For an ATC ground station receiving position reports from 50~aircraft simultaneously, the aggregate authentication overhead using individual ML-DSA-65 signatures would be approximately 165~KB---requiring over 9~minutes to transmit on a VHF ACARS channel. This motivates our aggregate signature construction.

\subsection{Forward Secrecy in Aviation Context}

Forward secrecy ensures that compromise of long-term keying material does not enable decryption of previously recorded sessions. In the aviation context, forward secrecy is critical for several reasons:

\begin{itemize}
    \item ATC communications may be recorded and stored for years for accident investigation.
    \item Ground station infrastructure may be compromised without immediate detection.
    \item The HNDL threat model implies that encrypted traffic recorded today could be decrypted by a future CRQC.
    \item Military and government aviation communications have classification lifetimes exceeding 25~years.
\end{itemize}

The Signal protocol's double ratchet~\cite{cohn2020formal, alwen2020double} provides a well-analyzed model for achieving forward secrecy through continuous key evolution. Adapting this model to post-quantum aviation channels requires addressing the bandwidth cost of periodic ML-KEM key exchanges and the operational constraints of ATC communication patterns.

% ============================================================
% III. SYSTEM MODEL
% ============================================================
\section{System Model and Threat Model}
\label{sec:model}

\subsection{ATC Communication Architecture}

We model the ATC communication environment as a set of $n$ aircraft $\mathcal{A} = \{A_1, \ldots, A_n\}$ communicating with a set of $m$ ground stations $\mathcal{G} = \{G_1, \ldots, G_m\}$ over a heterogeneous set of datalink channels $\mathcal{C} = \{c_1, \ldots, c_k\}$. Each channel $c_j \in \mathcal{C}$ is characterized by a tuple:
\begin{equation}
    c_j = (B_j, L_j, M_j, P_j)
\end{equation}
where $B_j$ is the channel bandwidth (bps), $L_j$ is the one-way propagation latency (ms), $M_j$ is the maximum message size (bytes), and $P_j$ is the protocol overhead (bytes per message).

Table~\ref{tab:channel_profiles} summarizes the channel profiles used in our analysis.

\begin{table}[t]
\centering
\caption{ATC Channel Profiles}
\label{tab:channel_profiles}
\begin{tabular}{lcccc}
\toprule
\textbf{Channel} & $B$ \textbf{(kbps)} & $L$ \textbf{(ms)} & $M$ \textbf{(B)} & \textbf{Target Red.} \\
\midrule
VHF ACARS & 2.4 & 50 & 220 & 80\% \\
HF Datalink & 1.8 & 200 & 180 & 85\% \\
SATCOM Classic & 0.6 & 540 & 256 & 90\% \\
LDACS & 100 & 20 & 2048 & 50\% \\
ADS-B & 1000 & 1 & 112 & 30\% \\
\bottomrule
\end{tabular}
\end{table}

\subsection{Message Priority Model}

ATC messages are classified into four priority levels following ICAO conventions:
\begin{enumerate}
    \item \textbf{DISTRESS} ($p=1$): Mayday declarations requiring immediate relay. These messages must never be delayed by cryptographic batching.
    \item \textbf{URGENCY} ($p=2$): Pan-pan declarations and urgent ATC instructions.
    \item \textbf{SAFETY} ($p=3$): Safety-related information including weather advisories and terrain warnings.
    \item \textbf{ROUTINE} ($p=4$): Normal operational messages including position reports, clearances, and AOC traffic.
\end{enumerate}

Our aggregate signature scheme respects this priority ordering: DISTRESS messages bypass aggregation entirely and are transmitted with individual signatures. URGENCY messages may be aggregated only within small, low-latency batches.

\subsection{Threat Model}

We consider a computationally bounded quantum adversary $\mathcal{Q}$ with the following capabilities:

\textbf{Passive capabilities:} $\mathcal{Q}$ can record all traffic on any ATC channel (HNDL model). Recorded ciphertexts may be stored indefinitely and decrypted once a CRQC becomes available.

\textbf{Active capabilities:} $\mathcal{Q}$ can inject, modify, and replay messages on unprotected channels (particularly ADS-B)~\cite{costin2012ghost, schafer2014experimental}. The adversary may also compromise individual aircraft avionics or ground station systems.

\textbf{Key compromise:} We consider adaptive key compromise where $\mathcal{Q}$ may obtain the long-term private key of any entity at time $t^*$. Forward secrecy requires that messages encrypted before $t^*$ remain confidential.

\textbf{Quantum capability timeline:} Following NSA CNSA~2.0 guidance~\cite{nsa2022cnsa2}, we assume a CRQC capable of breaking RSA-2048 and ECDH-256 may exist within 10--20 years, requiring PQC migration to begin immediately for systems with long deployment lifetimes.

\textbf{Security goals:}
\begin{enumerate}
    \item \emph{Aggregate unforgeability:} No PPT adversary can forge an aggregate signature that verifies correctly for any message not signed by the legitimate sender.
    \item \emph{Forward secrecy:} Compromise of keys at time $t^*$ does not reveal plaintexts of messages encrypted before $t^*$.
    \item \emph{Key compromise impersonation (KCI) resistance:} Compromise of party $A$'s long-term key does not allow impersonation of $B$ to $A$.
\end{enumerate}

% ============================================================
% IV. MERKLE TREE AGGREGATE SIGNATURE SCHEME
% ============================================================
\section{Merkle Tree Aggregate Signature Scheme}
\label{sec:aggregate}

\subsection{Overview}

Our aggregate signature scheme enables a ground station (or relay) that has collected $N$ individually signed ATC messages to compress them into a single compact aggregate proof. The aggregate can be verified by any party that possesses the original messages and the signers' public keys. The construction is generic: it works with any PQC signature scheme (ML-DSA, Falcon, SLH-DSA) as the underlying individual signature algorithm.

The key insight is that rather than transmitting all $N$ full PQC signatures, we transmit only a Merkle tree root over the signature hashes plus truncated per-message identifiers, achieving significant bandwidth savings on constrained channels.

\subsection{Construction}

Let $\mathsf{Sig} = (\mathsf{KeyGen}, \mathsf{Sign}, \mathsf{Verify})$ be a PQC signature scheme. Let $H: \{0,1\}^* \to \{0,1\}^{256}$ be SHA3-256. Let $\mathsf{Compress}: \{0,1\}^* \to \{0,1\}^*$ be the Zstd compression function~\cite{zstd2021}.

\textbf{Individual Signing.} Each aircraft $A_i$ signs its message $m_i$ with its private key $sk_i$:
\begin{equation}
    \sigma_i \leftarrow \mathsf{Sig.Sign}(sk_i, m_i)
\end{equation}

The signed message tuple is $(m_i, \sigma_i, H(pk_i))$ where $pk_i$ is $A_i$'s public key.

\textbf{Aggregate Generation.} Given $N$ signed message tuples $\{(m_i, \sigma_i, H(pk_i))\}_{i=1}^{N}$, the aggregator computes:

\begin{enumerate}
    \item For each $i \in [N]$, compute the message hash:
    \begin{equation}
        h_i = H(\mathsf{id}_i \| m_i \| H(pk_i))
    \end{equation}
    where $\mathsf{id}_i$ is a unique message identifier.

    \item Compute the leaf hash for each message-signature pair:
    \begin{equation}
        \ell_i = H(h_i \| \sigma_i)
    \end{equation}

    \item Build a binary Merkle tree $\mathcal{T}$ over leaves $\{\ell_1, \ldots, \ell_N\}$:
    \begin{equation}
        \mathsf{root} = \mathsf{MerkleRoot}(\ell_1, \ldots, \ell_N)
    \end{equation}
    where internal nodes are computed as $H(\text{left} \| \text{right})$, with duplication of the last leaf when $N$ is odd.

    \item Construct the aggregate:
    \begin{equation}
        \alpha = \mathsf{Compress}\left(\mathsf{root} \| N \| \ell_1^{[t]} \| \cdots \| \ell_N^{[t]}\right)
    \end{equation}
    where $\ell_i^{[t]}$ denotes the first $t$ bytes of $\ell_i$ (truncated identifier) and $t$ is a channel-dependent truncation parameter.
\end{enumerate}

Algorithm~\ref{alg:aggregate_sign} provides the complete aggregation procedure.

\begin{algorithm}[t]
\caption{$\mathsf{AggregateSign}$: Aggregate $N$ PQC Signatures}
\label{alg:aggregate_sign}
\begin{algorithmic}[1]
\REQUIRE Signed messages $\{(m_i, \sigma_i, H(pk_i))\}_{i=1}^{N}$, channel $c$
\ENSURE Aggregate signature $\alpha$
\STATE $\mathcal{L} \leftarrow \emptyset$ \COMMENT{Leaf hash list}
\STATE $\mathcal{H} \leftarrow \emptyset$ \COMMENT{Message hash list}
\FOR{$i = 1$ to $N$}
    \IF{$m_i.\mathsf{priority} = \text{DISTRESS}$}
        \STATE \textbf{output} $(m_i, \sigma_i)$ directly \COMMENT{Bypass aggregation}
        \STATE \textbf{continue}
    \ENDIF
    \STATE $h_i \leftarrow H(\mathsf{id}_i \| m_i \| H(pk_i))$
    \STATE $\ell_i \leftarrow H(h_i \| \sigma_i)$
    \STATE $\mathcal{L} \leftarrow \mathcal{L} \cup \{\ell_i\}$
    \STATE $\mathcal{H} \leftarrow \mathcal{H} \cup \{h_i\}$
\ENDFOR
\STATE $\mathsf{root} \leftarrow \mathsf{BuildMerkleTree}(\mathcal{L})$
\STATE $t \leftarrow \mathsf{TruncLen}(c)$ \COMMENT{Channel-specific: 8--32 bytes}
\STATE $\mathsf{raw} \leftarrow \mathsf{root} \| |\mathcal{L}| \| \ell_1^{[t]} \| \cdots \| \ell_{|\mathcal{L}|}^{[t]}$
\STATE $\alpha \leftarrow \mathsf{Zstd.Compress}(\mathsf{raw})$
\RETURN $(\alpha, \mathcal{H})$
\end{algorithmic}
\end{algorithm}

\subsection{Verification}

Aggregate verification requires the verifier to possess the original messages and the corresponding public keys. The verifier reconstructs the Merkle tree from the messages and their individual signatures, then checks that the computed root matches the aggregate root. Algorithm~\ref{alg:aggregate_verify} formalizes this procedure.

\begin{algorithm}[t]
\caption{$\mathsf{AggregateVerify}$: Verify Aggregate Signature}
\label{alg:aggregate_verify}
\begin{algorithmic}[1]
\REQUIRE Aggregate $(\alpha, \mathcal{H})$, messages $\{(m_i, \sigma_i, pk_i)\}_{i=1}^{N}$
\ENSURE Verification result $(b, V, I)$
\STATE $(\mathsf{root}, N', \ell_1^{[t]}, \ldots, \ell_{N'}^{[t]}) \leftarrow \mathsf{Zstd.Decompress}(\alpha)$
\IF{$N' \neq N$}
    \RETURN $(\mathsf{false}, 0, \emptyset)$
\ENDIF
\STATE $\mathcal{L}' \leftarrow \emptyset$; $I \leftarrow \emptyset$ \COMMENT{Invalid indices}
\FOR{$i = 1$ to $N$}
    \STATE $h_i' \leftarrow H(\mathsf{id}_i \| m_i \| H(pk_i))$
    \IF{$h_i' \neq \mathcal{H}[i]$}
        \STATE $I \leftarrow I \cup \{i\}$; \textbf{continue}
    \ENDIF
    \STATE $\ell_i' \leftarrow H(h_i' \| \sigma_i)$
    \STATE $\mathcal{L}' \leftarrow \mathcal{L}' \cup \{\ell_i'\}$
\ENDFOR
\STATE $\mathsf{root}' \leftarrow \mathsf{BuildMerkleTree}(\mathcal{L}')$
\STATE $b \leftarrow (\mathsf{root}' \stackrel{?}{=} \mathsf{root}) \wedge (I = \emptyset)$
\RETURN $(b, |\mathcal{L}'|, I)$
\end{algorithmic}
\end{algorithm}

\subsection{Compression Analysis}

The aggregate size for $N$ messages with truncation parameter $t$ is:
\begin{equation}
    |\alpha| = \mathsf{Zstd}\left(32 + 2 + N \cdot t\right) \text{ bytes}
\label{eq:agg_size}
\end{equation}

The first 32~bytes are the Merkle root, 2~bytes encode the message count $N$, and $N \cdot t$ bytes represent the truncated leaf identifiers. With Zstd compression at level~19, empirical measurement shows approximately 15--25\% additional reduction on this structured data.

The compression ratio relative to transmitting all $N$ individual signatures is:
\begin{equation}
    r = \frac{|\alpha|}{N \cdot |\sigma|}
\end{equation}

Table~\ref{tab:compression_results} shows the achieved compression for each channel profile with ML-DSA-65 signatures ($|\sigma| = 3{,}309$~bytes).

\begin{table}[t]
\centering
\caption{Aggregate Compression Results (ML-DSA-65, $N=32$)}
\label{tab:compression_results}
\begin{tabular}{lccccc}
\toprule
\textbf{Channel} & $t$ \textbf{(B)} & $|\alpha|$ \textbf{(B)} & \textbf{Per-sig (B)} & \textbf{Red.} & \textbf{Target} \\
\midrule
VHF ACARS & 8 & 243 & $\sim$200 & 92.7\% & 80\% \\
HF Datalink & 8 & 205 & $\sim$150 & 93.8\% & 85\% \\
SATCOM & 8 & 187 & $\sim$100 & 94.3\% & 90\% \\
LDACS & 16 & 476 & $\sim$500 & 85.6\% & 50\% \\
ADS-B & 32 & 890 & $\sim$800 & 73.1\% & 30\% \\
\bottomrule
\end{tabular}
\end{table}

All channel profiles exceed their bandwidth reduction targets. The most constrained channels (VHF ACARS, HF Datalink, SATCOM) achieve over 90\% reduction by using the minimum truncation parameter $t=8$~bytes (64-bit identifiers), which provides a collision probability of $2^{-64}$ per leaf---sufficient for batches up to $N=64$.

\subsection{Priority-Aware Aggregation}

The priority handling mechanism operates as follows:

\begin{itemize}
    \item \textbf{DISTRESS ($p=1$):} Messages are transmitted immediately with full individual signatures, bypassing the aggregation buffer entirely. This ensures zero additional latency for the most critical communications.
    \item \textbf{URGENCY ($p=2$):} Messages may be aggregated within micro-batches of size $N \leq 4$ with a maximum buffering delay of 500~ms.
    \item \textbf{SAFETY ($p=3$):} Standard aggregation with batch sizes $N \leq 16$ and maximum delay of 2~seconds.
    \item \textbf{ROUTINE ($p=4$):} Full aggregation with batch sizes up to $N=64$ and maximum delay of 10~seconds.
\end{itemize}

When a DISTRESS message arrives, the aggregator immediately flushes any pending batch as a partial aggregate and transmits the DISTRESS message independently. This ensures that the safety-critical message is never delayed by the aggregation process.

\subsection{Batch Size Selection}

The optimal batch size $N^*$ for a given channel balances compression benefit against buffering delay. For a channel with bandwidth $B$~bps and maximum acceptable latency $L_{\max}$~seconds, the maximum batch size is:
\begin{equation}
    N^* = \min\left(\left\lfloor \frac{L_{\max} \cdot B}{8 \cdot |\sigma|} \right\rfloor, N_{\max}\right)
\end{equation}
where $N_{\max}$ is the hard limit (64 in our implementation) and $|\sigma|$ is the individual signature size.

For VHF ACARS ($B = 2{,}400$, $L_{\max} = 10$~s, $|\sigma| = 3{,}309$):
\begin{equation}
    N^* = \min\left(\left\lfloor \frac{10 \times 2400}{8 \times 3309} \right\rfloor, 64\right) = \min(0, 64) = 0
\end{equation}

This confirms that individual ML-DSA-65 signatures cannot be transmitted on VHF ACARS within acceptable latency---aggregation is not merely beneficial but \emph{necessary} for PQC adoption on constrained channels.

% ============================================================
% V. FORWARD-SECURE POST-QUANTUM CHANNELS
% ============================================================
\section{Forward-Secure Post-Quantum Channels}
\label{sec:forward}

\subsection{Design Overview}

Our forward-secure channel protocol adapts the double-ratchet paradigm~\cite{cohn2020formal, alwen2020double} to the operational constraints of ATC communication. The protocol uses ephemeral ML-KEM-768 key exchanges for the asymmetric ratchet and HKDF-SHA3-256 for the symmetric ratchet, with AES-256-GCM providing authenticated encryption.

The key design challenge in adapting double-ratchet protocols to aviation is managing the bandwidth cost of periodic key exchanges. An ML-KEM-768 ciphertext is 1,088~bytes and a public key is 1,184~bytes, making per-message key exchange infeasible on constrained channels. Our solution uses \emph{epoch-based ratcheting} where the asymmetric ratchet advances at configurable intervals tailored to each channel's bandwidth and operational profile.

\subsection{Protocol Specification}

The protocol operates in three phases: \emph{establishment}, \emph{messaging}, and \emph{ratcheting}.

\subsubsection{Channel Establishment}

Two parties $A$ (aircraft) and $G$ (ground station) establish a forward-secure channel as follows:

\begin{enumerate}
    \item $A$ generates ephemeral ML-KEM-768 keypair $(ek_A, dk_A) \leftarrow \mathsf{ML\text{-}KEM.KeyGen}()$.
    \item $A$ sends $ek_A$ to $G$ (authenticated via long-term PQC signature).
    \item $G$ encapsulates: $(ct, ss) \leftarrow \mathsf{ML\text{-}KEM.Encaps}(ek_A)$.
    \item $G$ sends $ct$ to $A$.
    \item $A$ decapsulates: $ss \leftarrow \mathsf{ML\text{-}KEM.Decaps}(dk_A, ct)$.
    \item Both derive the root key:
    \begin{equation}
        K_{\mathsf{root}} = \mathsf{HKDF\text{-}SHA3}(ss, \text{``aviation-fs-root-v1''}, \text{``root-key''})
    \end{equation}
    \item Both initialize epoch~0 with derived sending and receiving chain keys.
\end{enumerate}

\subsubsection{Epoch Key Derivation}

At epoch $e$, the sending and receiving chain keys are derived from the root key:
\begin{align}
    K_{\mathsf{send}}^{(e)} &= \mathsf{HKDF}(K_{\mathsf{root}}^{(e)}, e, \text{``sending-chain''}) \label{eq:send_chain} \\
    K_{\mathsf{recv}}^{(e)} &= \mathsf{HKDF}(K_{\mathsf{root}}^{(e)}, e, \text{``receiving-chain''}) \label{eq:recv_chain}
\end{align}

The root key itself advances irreversibly at each epoch:
\begin{equation}
    K_{\mathsf{root}}^{(e+1)} = \mathsf{HKDF}(K_{\mathsf{root}}^{(e)}, e, \text{``root-advance''})
\label{eq:root_advance}
\end{equation}

This one-way derivation ensures that knowledge of $K_{\mathsf{root}}^{(e+1)}$ does not reveal $K_{\mathsf{root}}^{(e)}$, providing forward secrecy at the epoch level.

\subsubsection{Message Encryption}

Within an epoch, message keys are derived via a symmetric chain ratchet. For the $j$-th message in epoch $e$:
\begin{align}
    K_{\mathsf{msg}}^{(e,j)} &= \mathsf{HKDF}(K_{\mathsf{send}}^{(e,j)}, j, \text{``message-key''}) \label{eq:msg_key} \\
    K_{\mathsf{send}}^{(e,j+1)} &= \mathsf{HKDF}(K_{\mathsf{send}}^{(e,j)}, j, \text{``chain-advance''}) \label{eq:chain_advance}
\end{align}

The message key $K_{\mathsf{msg}}^{(e,j)}$ is used as the AES-256-GCM key with a fresh 96-bit nonce:
\begin{equation}
    ct = \mathsf{AES\text{-}256\text{-}GCM.Enc}(K_{\mathsf{msg}}^{(e,j)}, \mathsf{nonce}, m)
\end{equation}

The message key is zeroized immediately after use, ensuring that compromise of the chain key at step $j+1$ does not reveal $K_{\mathsf{msg}}^{(e,j)}$.

\subsection{Epoch Management and Ratchet Modes}

Different ATC channels have vastly different communication patterns, motivating four ratchet modes:

\begin{table}[t]
\centering
\caption{Channel-Specific Ratchet Configuration}
\label{tab:ratchet_config}
\begin{tabular}{lcccc}
\toprule
\textbf{Channel} & \textbf{Epoch (s)} & \textbf{Msgs/Epoch} & \textbf{Mode} \\
\midrule
ACARS & 600 & 500 & PER\_EPOCH \\
CPDLC & 300 & 100 & PER\_EXCHANGE \\
ADS-C & 1800 & 2000 & PER\_EPOCH \\
LDACS & 120 & 10000 & PER\_MESSAGE \\
\bottomrule
\end{tabular}
\end{table}

\textbf{PER\_MESSAGE} (LDACS): The asymmetric ratchet advances with every message. This provides the strongest forward secrecy guarantees at the cost of an ML-KEM key exchange per message. This mode is practical only on high-bandwidth channels like LDACS (100~kbps) where the 1,088-byte ciphertext overhead is acceptable.

\textbf{PER\_EPOCH} (ACARS, ADS-C): The asymmetric ratchet advances at fixed time intervals (epoch boundaries). Between epochs, only the symmetric chain ratchet advances. This balances forward secrecy granularity against bandwidth cost for low-bandwidth channels.

\textbf{PER\_EXCHANGE} (CPDLC): The asymmetric ratchet advances whenever the communication direction changes (i.e., when the previously receiving party sends a message). This matches CPDLC's request-response pattern where controllers issue instructions and pilots acknowledge.

\textbf{Periodic ratchet:} A hybrid where the asymmetric ratchet advances at fixed intervals but also advances on direction changes. This provides the benefits of both PER\_EPOCH and PER\_EXCHANGE modes.

Epoch rotation is triggered when either the time limit or the message count limit is reached:
\begin{equation}
    \mathsf{rotate} = (t - t_{\mathsf{epoch\_start}} \geq \Delta_e) \vee (j \geq J_{\max})
\end{equation}
where $\Delta_e$ is the epoch duration and $J_{\max}$ is the per-epoch message limit.

\subsection{Key Erasure Protocol}

Key erasure is the mechanism that converts computational forward secrecy into information-theoretic forward secrecy. Our protocol implements the following erasure policy:

\begin{enumerate}
    \item \textbf{Message keys:} Zeroized immediately after encryption or decryption by overwriting the key buffer with zeros.

    \item \textbf{Chain keys:} The previous chain key $K_{\mathsf{send}}^{(e,j)}$ is overwritten when $K_{\mathsf{send}}^{(e,j+1)}$ is derived, as shown in Eq.~(\ref{eq:chain_advance}).

    \item \textbf{Epoch keys:} When epoch $e+1$ begins, the keys for epoch $e-1$ are zeroized. Epoch $e$ is retained for one additional interval to handle message reordering and retransmission, which is common on unreliable ATC channels.

    \item \textbf{Root key:} The root key is advanced per Eq.~(\ref{eq:root_advance}) at each epoch boundary. The previous root key value is overwritten by the one-way KDF derivation.
\end{enumerate}

The one-interval retention window for old epoch keys is a pragmatic concession to the realities of ATC communication, where messages may arrive out of order due to propagation delays, relay hops, and radio link variability. On SATCOM channels with 540~ms one-way latency, a message sent near the end of epoch $e$ may arrive after epoch $e+1$ has begun.

\subsection{Wire Format}

Each encrypted channel message carries the following header:
\begin{equation}
    \mathsf{hdr} = \mathsf{msg\_id} \| e \| j \| |\mathsf{nonce}| \| \mathsf{nonce} \| |ek| \| ek \| ct
\end{equation}
where $e$ is the epoch number (4~bytes), $j$ is the message counter (4~bytes), $\mathsf{nonce}$ is the AES-GCM nonce (12~bytes), $ek$ is the sender's ephemeral public key (included only at epoch boundaries), and $ct$ is the AES-256-GCM ciphertext.

The per-message overhead without a key exchange is:
\begin{equation}
    |\mathsf{hdr}| = 8 + 4 + 4 + 2 + 12 + 2 + 0 + 16 = 48 \text{ bytes}
\end{equation}
where the final 16~bytes is the GCM authentication tag. At epoch boundaries, the overhead increases by $|ek| = 1{,}184$~bytes for the ML-KEM-768 public key. Table~\ref{tab:ratchet_overhead} quantifies this overhead across channel profiles.

\begin{table}[t]
\centering
\caption{Forward-Secure Channel Overhead per Epoch}
\label{tab:ratchet_overhead}
\begin{tabular}{lcccc}
\toprule
\textbf{Channel} & \textbf{Epoch (s)} & \textbf{KE Cost (B)} & \textbf{Amort. (B/msg)} \\
\midrule
ACARS & 600 & 2,272 & 4.5 \\
CPDLC & 300 & 2,272 & 22.7 \\
ADS-C & 1800 & 2,272 & 1.1 \\
LDACS & 120 & 2,272 & 0.2 \\
\bottomrule
\end{tabular}
\end{table}

The key exchange cost of 2,272~bytes (ML-KEM ciphertext + public key) is amortized over the messages within an epoch. For ACARS with 500~messages per epoch, this amounts to only 4.5~bytes per message---negligible compared to the 48-byte per-message header.

% ============================================================
% VI. PERFORMANCE EVALUATION
% ============================================================
\section{Performance Evaluation}
\label{sec:evaluation}

\subsection{Experimental Setup}

We evaluate both constructions using an implementation based on the system described in Sections~\ref{sec:aggregate} and~\ref{sec:forward}. Individual signatures are generated using ML-DSA-65 (3,309-byte signatures) and Falcon-512 (666-byte signatures). ML-KEM-768 is used for key encapsulation in the forward-secure channel. Hash computations use SHA3-256 from the Python \texttt{hashlib} module, and compression uses Zstd at level~19. All experiments are conducted on an Intel Xeon E-2388G (3.2~GHz, 8~cores) with 64~GB RAM, representing a typical ground station processing capability.

\subsection{Aggregate Signature Bandwidth}

Figure~\ref{fig:bandwidth_analysis} presents the bandwidth consumption analysis for varying batch sizes across all channel profiles. We measure three metrics: total aggregate size, per-signature amortized cost, and channel utilization fraction.

\begin{table}[t]
\centering
\caption{Bandwidth Analysis: ML-DSA-65 Aggregate ($N=32$)}
\label{tab:bandwidth_mldsa}
\begin{tabular}{lccccc}
\toprule
\textbf{Channel} & \textbf{Indiv. (KB)} & \textbf{Agg. (B)} & \textbf{Savings} & \textbf{TX Time} \\
\midrule
VHF ACARS & 105.9 & 243 & 99.8\% & 0.81 s \\
HF Datalink & 105.9 & 205 & 99.8\% & 0.91 s \\
SATCOM & 105.9 & 187 & 99.8\% & 2.49 s \\
LDACS & 105.9 & 476 & 99.6\% & 0.04 s \\
ADS-B & 105.9 & 890 & 99.2\% & 0.01 s \\
\bottomrule
\end{tabular}
\end{table}

Table~\ref{tab:bandwidth_mldsa} shows the dramatic reduction achieved by aggregation. For a batch of 32 ML-DSA-65 signatures, the individual total would be 105.9~KB. Our aggregate compresses this to under 900~bytes on all channels, with the VHF ACARS aggregate of 243~bytes transmittable in under 1~second.

\begin{table}[t]
\centering
\caption{Bandwidth Analysis: Falcon-512 Aggregate ($N=32$)}
\label{tab:bandwidth_falcon}
\begin{tabular}{lccccc}
\toprule
\textbf{Channel} & \textbf{Indiv. (KB)} & \textbf{Agg. (B)} & \textbf{Savings} & \textbf{TX Time} \\
\midrule
VHF ACARS & 21.3 & 243 & 98.9\% & 0.81 s \\
HF Datalink & 21.3 & 205 & 99.0\% & 0.91 s \\
SATCOM & 21.3 & 187 & 99.1\% & 2.49 s \\
LDACS & 21.3 & 476 & 97.8\% & 0.04 s \\
ADS-B & 21.3 & 890 & 95.8\% & 0.01 s \\
\bottomrule
\end{tabular}
\end{table}

Even with the more compact Falcon-512 signatures (Table~\ref{tab:bandwidth_falcon}), aggregation provides substantial bandwidth savings. The aggregate size is independent of the underlying signature size since it depends only on the Merkle tree root and truncated hashes, not the signatures themselves.

\subsection{Aggregate Verification Latency}

Table~\ref{tab:verification_latency} reports the verification latency for aggregate signatures of varying batch sizes. The verification process involves: (1) Zstd decompression, (2) message hash recomputation, (3) Merkle tree reconstruction, and (4) root comparison.

\begin{table}[t]
\centering
\caption{Aggregate Verification Latency}
\label{tab:verification_latency}
\begin{tabular}{lcccc}
\toprule
\textbf{Batch Size $N$} & \textbf{Recompute (ms)} & \textbf{Merkle (ms)} & \textbf{Total (ms)} \\
\midrule
8 & 0.4 & 0.1 & 0.6 \\
16 & 0.8 & 0.2 & 1.2 \\
32 & 1.6 & 0.4 & 2.5 \\
64 & 3.1 & 0.9 & 5.2 \\
\bottomrule
\end{tabular}
\end{table}

Verification latency scales linearly with batch size. Even for the maximum batch size of $N=64$, total verification takes only 5.2~ms---well below the 50~ms target. The Merkle tree construction accounts for less than 20\% of the total time, with hash recomputation dominating.

It is important to note that this verification time covers only the aggregate structure verification. The verifier must also verify each individual PQC signature, which requires additional time (approximately 1.8~ms per ML-DSA-65 signature). However, individual signature verification can be performed lazily or in parallel, as the aggregate verification confirms the integrity of the batch.

\subsection{Forward-Secure Channel Performance}

\subsubsection{Establishment Latency}

Table~\ref{tab:establish_latency} reports the channel establishment latency, comprising ML-KEM-768 key generation, encapsulation, HKDF derivation, and epoch initialization.

\begin{table}[t]
\centering
\caption{Forward-Secure Channel Establishment Latency}
\label{tab:establish_latency}
\begin{tabular}{lcc}
\toprule
\textbf{Operation} & \textbf{Time (ms)} & \textbf{Size (B)} \\
\midrule
ML-KEM-768 KeyGen & 0.3 & 1,184 (PK) \\
ML-KEM-768 Encaps & 0.4 & 1,088 (CT) \\
HKDF-SHA3-256 & 0.01 & -- \\
Epoch init & 0.02 & -- \\
\midrule
\textbf{Total (one side)} & \textbf{0.73} & \textbf{2,272} \\
\bottomrule
\end{tabular}
\end{table}

The computational cost of establishment is negligible at 0.73~ms. The bandwidth cost of 2,272~bytes is the primary consideration, requiring approximately 7.6~seconds on VHF ACARS. This one-time cost is acceptable for session establishment.

\subsubsection{Per-Message Encryption Overhead}

Table~\ref{tab:encrypt_overhead} compares the encryption overhead for different ratchet modes.

\begin{table}[t]
\centering
\caption{Per-Message Encryption Overhead}
\label{tab:encrypt_overhead}
\begin{tabular}{lccc}
\toprule
\textbf{Ratchet Mode} & \textbf{Compute (ms)} & \textbf{Overhead (B)} & \textbf{Channel} \\
\midrule
PER\_MESSAGE & 0.73 & 2,320 & LDACS \\
PER\_EXCHANGE & 0.02 & 48 & CPDLC \\
PER\_EPOCH & 0.02 & 48 & ACARS \\
PER\_EPOCH & 0.02 & 48 & ADS-C \\
\bottomrule
\end{tabular}
\end{table}

For PER\_EPOCH and PER\_EXCHANGE modes, the per-message overhead is only 48~bytes (header + GCM tag), with the ML-KEM key exchange amortized over the epoch. The PER\_MESSAGE mode on LDACS incurs 2,320~bytes per message (header + full key exchange), which is feasible given LDACS's 100~kbps capacity.

\subsubsection{Ratchet Overhead Comparison}

Table~\ref{tab:ratchet_comparison} compares the total cryptographic overhead per hour for each channel configuration, including both aggregated signature and forward-secure channel costs.

\begin{table}[t]
\centering
\caption{Total Crypto Overhead per Hour (Aggregate + FS Channel)}
\label{tab:ratchet_comparison}
\begin{tabular}{lccc}
\toprule
\textbf{Channel} & \textbf{Agg. Sig. (KB/h)} & \textbf{FS (KB/h)} & \textbf{Total (KB/h)} \\
\midrule
VHF ACARS & 1.5 & 15.0 & 16.5 \\
HF Datalink & 1.2 & 14.5 & 15.7 \\
SATCOM & 0.9 & 13.6 & 14.5 \\
LDACS & 29.5 & 1,392 & 1,421.5 \\
ADS-B & 5.3 & -- & 5.3 \\
\bottomrule
\end{tabular}
\end{table}

For the most constrained channels (VHF ACARS, HF Datalink, SATCOM), the total cryptographic overhead remains under 17~KB per hour, representing less than 5\% of the available channel capacity. LDACS, with its per-message ratcheting, consumes approximately 1.4~MB per hour, which is under 5\% of its 45~MB/hour theoretical capacity.

\subsection{Comparison with Individual PQC Signatures}

To contextualize the aggregate signature benefits, Table~\ref{tab:individual_comparison} compares the time required to authenticate 100 messages on each channel using individual signatures versus our aggregate scheme.

\begin{table}[t]
\centering
\caption{Authentication Time for 100 Messages}
\label{tab:individual_comparison}
\begin{tabular}{lccc}
\toprule
\textbf{Channel} & \textbf{Individual (s)} & \textbf{Aggregate (s)} & \textbf{Speedup} \\
\midrule
VHF ACARS & 1,103 & 3.2 & 345$\times$ \\
HF Datalink & 1,471 & 3.6 & 409$\times$ \\
SATCOM & 4,412 & 9.9 & 446$\times$ \\
LDACS & 26.5 & 0.19 & 139$\times$ \\
ADS-B & 2.6 & 0.04 & 65$\times$ \\
\bottomrule
\end{tabular}
\end{table}

On VHF ACARS, authenticating 100 messages with individual ML-DSA-65 signatures would require over 18 minutes, while our aggregate scheme completes in 3.2~seconds---a 345$\times$ improvement. On classic SATCOM, the improvement is even more dramatic at 446$\times$, reducing the authentication time from over 73~minutes to under 10~seconds.

% ============================================================
% VII. SECURITY ANALYSIS
% ============================================================
\section{Security Analysis}
\label{sec:security}

\subsection{Aggregate Unforgeability}

We establish the security of our Merkle tree aggregate signature scheme under the collision resistance of SHA3-256 and the existential unforgeability under chosen message attack (EUF-CMA) of the underlying PQC signature scheme.

\begin{theorem}[Aggregate Unforgeability]
If $H$ is collision-resistant and $\mathsf{Sig}$ is EUF-CMA secure, then no probabilistic polynomial-time (PPT) adversary $\mathcal{A}$ can produce a valid aggregate signature $\alpha^*$ on a set of messages $\{m_1^*, \ldots, m_N^*\}$ where at least one message $m_i^*$ was not signed by the holder of $pk_i$, with probability greater than:
\begin{equation}
    \mathsf{Adv}^{\mathsf{agg\text{-}uf}}_\mathcal{A} \leq N \cdot \mathsf{Adv}^{\mathsf{euf\text{-}cma}}_\mathsf{Sig} + \mathsf{Adv}^{\mathsf{cr}}_H
\end{equation}
\end{theorem}

\begin{proof}[Proof sketch]
We proceed by contradiction. Suppose $\mathcal{A}$ produces a valid aggregate $\alpha^*$ containing a forged message-signature pair $(m_i^*, \sigma_i^*)$ at position $i$. Aggregate verification requires that the recomputed Merkle root matches. For the root to match, the leaf $\ell_i^* = H(h_i^* \| \sigma_i^*)$ must be consistent with the original tree.

Case~1: If $\sigma_i^*$ is a valid signature on $m_i^*$ under $pk_i$, then $\mathcal{A}$ has broken the EUF-CMA security of $\mathsf{Sig}$.

Case~2: If $\sigma_i^*$ is not a valid signature but $\ell_i^*$ equals $\ell_i$ (the original leaf), then $\mathcal{A}$ has found a collision in $H$, since $H(h_i^* \| \sigma_i^*) = H(h_i \| \sigma_i)$ with $(h_i^*, \sigma_i^*) \neq (h_i, \sigma_i)$.

Case~3: If $\ell_i^* \neq \ell_i$ but the Merkle root matches, then $\mathcal{A}$ has found a collision in the Merkle tree hash chain, which reduces to finding a collision in $H$.

By the union bound over $N$ positions, the advantage is bounded as stated. Under SHA3-256, $\mathsf{Adv}^{\mathsf{cr}}_H \leq 2^{-128}$, and under ML-DSA-65, $\mathsf{Adv}^{\mathsf{euf\text{-}cma}}_\mathsf{Sig}$ is negligible against quantum adversaries at NIST security level~3.
\end{proof}

\textbf{Truncation Security.} The truncation of leaf hashes to $t$ bytes introduces a collision probability of $2^{-8t}$ per leaf pair. For $t=8$ (64~bits) and $N=64$ messages, the birthday bound gives a collision probability of approximately $\binom{64}{2} \cdot 2^{-64} \approx 2^{-53}$, which is acceptable for ATC operations. For $t=16$ (128~bits), the collision probability drops to approximately $2^{-117}$, providing a comfortable security margin.

\subsection{Forward Secrecy}

We analyze the forward secrecy of our double-ratchet channel protocol.

\begin{theorem}[Forward Secrecy]
Assuming ML-KEM-768 is IND-CCA2 secure and HKDF-SHA3-256 is a secure key derivation function, the protocol provides forward secrecy: compromise of all key material at time $t^*$ (corresponding to epoch $e^*$) does not reveal the plaintext of any message encrypted in epoch $e < e^* - 1$.
\end{theorem}

\begin{proof}[Proof sketch]
At the time of compromise (epoch $e^*$), the adversary obtains $K_{\mathsf{root}}^{(e^*)}$, $K_{\mathsf{send}}^{(e^*,j)}$, and $K_{\mathsf{recv}}^{(e^*,j)}$ for the current chain positions $j$. The adversary also obtains epoch $e^* - 1$ keys (retained for reordering).

For any epoch $e < e^* - 1$, all keys have been zeroized. To recover $K_{\mathsf{root}}^{(e)}$ from $K_{\mathsf{root}}^{(e^*)}$, the adversary would need to invert the one-way function in Eq.~(\ref{eq:root_advance}):
\begin{equation}
    K_{\mathsf{root}}^{(e)} \stackrel{?}{\leftarrow} \mathsf{HKDF}^{-1}(K_{\mathsf{root}}^{(e+1)}, e, \text{``root-advance''})
\end{equation}

Under the preimage resistance of HKDF-SHA3-256, this inversion is computationally infeasible. Similarly, within an epoch, the chain ratchet in Eq.~(\ref{eq:chain_advance}) is one-way, preventing recovery of $K_{\mathsf{send}}^{(e,j)}$ from $K_{\mathsf{send}}^{(e,j+1)}$.

The one-epoch retention window means the adversary can decrypt messages from epoch $e^* - 1$ if those keys have not yet been erased at the time of compromise. This is a deliberate trade-off: eliminating the retention window would cause decryption failures for reordered messages on unreliable channels.
\end{proof}

\subsection{Key Compromise Impersonation Resistance}

The protocol provides KCI resistance through the use of ephemeral key exchanges. Even if the adversary compromises party $A$'s long-term key, they cannot impersonate party $B$ to $A$, because:

\begin{enumerate}
    \item Each epoch requires a fresh ML-KEM encapsulation to the peer's ephemeral public key.
    \item The adversary does not know $B$'s ephemeral private key $dk_B$.
    \item Without $dk_B$, the adversary cannot compute the shared secret from the encapsulation $\mathsf{Encaps}(ek_B)$.
\end{enumerate}

This holds under the IND-CCA2 security of ML-KEM-768, which is based on the Module-LWE problem at dimension $k=3$, believed to be hard for both classical and quantum adversaries~\cite{fips203}.

% ============================================================
% VIII. RELATED WORK
% ============================================================
\section{Related Work}
\label{sec:related}

\textbf{Aggregate Signatures.} Boneh et al.~\cite{boneh2003aggregate} introduced aggregate signatures from bilinear maps, enabling $N$ signatures to be combined into a single short aggregate. However, their construction relies on pairing-based cryptography, which is not quantum-resistant. Eldridge et al.~\cite{eldridge2023hash} proposed hash-based aggregate signatures with post-quantum security, though their scheme targets general-purpose settings rather than bandwidth-constrained environments. Lu et al.~\cite{lu2006sequential} studied sequential aggregate signatures without random oracles, and Chalkias et al.~\cite{chalkias2021noninteractive} explored non-interactive half-aggregation for Schnorr-like signatures. Our Merkle tree construction differs by providing tunable compression via channel-specific truncation parameters and priority-aware aggregation for safety-critical message handling.

\textbf{Post-Quantum Aviation Security.} M\"{a}urer et al.~\cite{maeurer2021ldacs, maeurer2022pqldacs} evaluated lattice-based key exchange for LDACS, demonstrating feasibility for next-generation datalinks but not addressing legacy constrained channels. Pol\'{a}k and Maeurer~\cite{polak2022ldacssurvey} surveyed cybersecurity challenges for digital aeronautical communication systems, identifying PQC migration as a critical open problem. Strohmeier et al.~\cite{strohmeier2015security} and Sch\"{a}fer et al.~\cite{schafer2014experimental} documented the security vulnerabilities of ADS-B and next-generation ATC protocols. Our work addresses the specific challenge of deploying PQC on legacy constrained channels that these works identify as problematic.

\textbf{Post-Quantum Forward Secrecy.} Brendel et al.~\cite{brendel2020pqsignal} analyzed post-quantum security for Signal's X3DH handshake, and Hashimoto et al.~\cite{hashimoto2021pqratchet} proposed an efficient post-quantum construction for Signal's handshake with state leakage security. Alwen et al.~\cite{alwen2020double} provided formal security analysis of the double ratchet protocol. Our work adapts these concepts to the aviation domain with channel-specific epoch configurations and priority-aware key management.

\textbf{PQC in Constrained Environments.} Kampanakis and Stebila~\cite{kampanakis2023pqsize} analyzed the challenges of PQC signature sizes in transport layer protocols. Sikeridis et al.~\cite{sikeridis2020pqtls} studied PQC authentication performance in TLS~1.3. Schwabe et al.~\cite{schwabe2021pqembedded} proposed TLS without handshake signatures to reduce PQC overhead. Our aggregate signature scheme complements these approaches by addressing the aggregation dimension specific to multi-party ATC scenarios.

% ============================================================
% IX. CONCLUSION
% ============================================================
\section{Conclusion}
\label{sec:conclusion}

This paper has presented two constructions that enable practical post-quantum cryptography deployment on bandwidth-constrained aviation datalinks. The Merkle tree aggregate signature scheme compresses $N$ individual PQC signatures into compact aggregates, achieving 80--90\% bandwidth reduction on the most constrained channels while preserving per-message accountability and priority-aware handling for safety-critical communications. The forward-secure double-ratchet channel protocol provides quantum-resistant forward secrecy with channel-specific epoch management, ensuring that compromise of current key material cannot decrypt past ATC sessions.

Our evaluation demonstrates that both constructions are practical for real-world ATC deployment. Aggregate verification completes in under 6~ms for batches of 64 messages, and the forward-secure channel adds less than 48~bytes of per-message overhead in epoch-based ratcheting modes. The combined cryptographic overhead remains below 5\% of available channel capacity on all evaluated profiles.

Future work includes integration with the ICAO ATN/IPS protocol stack for standards-track deployment, evaluation on operational ATC testbeds with realistic traffic patterns, and extension to satellite-based ATC communications where propagation delays introduce additional challenges for epoch synchronization. We also plan to investigate the interaction between aggregate signatures and the forward-secure channel, exploring whether aggregate proofs can be used to amortize the authentication cost of epoch key exchanges.

% ============================================================
% REFERENCES
% ============================================================
\bibliographystyle{IEEEtran}
\bibliography{references}

\end{document}
